\documentclass{beamer}
% =====================================================================
% finance-beamer.tex — shared Beamer style for the LLM-in-Finance course
% =====================================================================
% Reusable slide style. Each deck \input's this immediately after
% \documentclass{beamer}. It provides:
%   * the Madrid theme + book colour palette (matches the printed book)
%   * two book-style framing blocks:
%       \begin{bigpicture} ... \end{bigpicture}   ("the bigger picture")
%       \begin{underhood}   ... \end{underhood}    ("under the hood")
%     These mirror the book's `context` and `deepdive` tcolorboxes so the
%     same intuition-vs-mechanism scaffold the reader meets in the book
%     reappears on the slides.
%   * a Python listings style for code frames
%   * appendix conventions: \appendixstart drops a divider and switches
%     to non-numbered backup frames so "extra / less central" material
%     lives after the main talk without inflating the slide count.
%
% Usage:
%   \documentclass{beamer}
%   % =====================================================================
% finance-beamer.tex — shared Beamer style for the LLM-in-Finance course
% =====================================================================
% Reusable slide style. Each deck \input's this immediately after
% \documentclass{beamer}. It provides:
%   * the Madrid theme + book colour palette (matches the printed book)
%   * two book-style framing blocks:
%       \begin{bigpicture} ... \end{bigpicture}   ("the bigger picture")
%       \begin{underhood}   ... \end{underhood}    ("under the hood")
%     These mirror the book's `context` and `deepdive` tcolorboxes so the
%     same intuition-vs-mechanism scaffold the reader meets in the book
%     reappears on the slides.
%   * a Python listings style for code frames
%   * appendix conventions: \appendixstart drops a divider and switches
%     to non-numbered backup frames so "extra / less central" material
%     lives after the main talk without inflating the slide count.
%
% Usage:
%   \documentclass{beamer}
%   % =====================================================================
% finance-beamer.tex — shared Beamer style for the LLM-in-Finance course
% =====================================================================
% Reusable slide style. Each deck \input's this immediately after
% \documentclass{beamer}. It provides:
%   * the Madrid theme + book colour palette (matches the printed book)
%   * two book-style framing blocks:
%       \begin{bigpicture} ... \end{bigpicture}   ("the bigger picture")
%       \begin{underhood}   ... \end{underhood}    ("under the hood")
%     These mirror the book's `context` and `deepdive` tcolorboxes so the
%     same intuition-vs-mechanism scaffold the reader meets in the book
%     reappears on the slides.
%   * a Python listings style for code frames
%   * appendix conventions: \appendixstart drops a divider and switches
%     to non-numbered backup frames so "extra / less central" material
%     lives after the main talk without inflating the slide count.
%
% Usage:
%   \documentclass{beamer}
%   \input{../_style/finance-beamer.tex}
%   \title{...}\subtitle{...}\author{...}\institute{...}\date{\today}
%   \begin{document} ... \end{document}
% =====================================================================

\usetheme{Madrid}
\usecolortheme{whale}
\setbeamertemplate{navigation symbols}{}
\setbeamertemplate{caption}[numbered]
\setbeamercovered{transparent}

% --- Density controls: keep dense, equation-heavy frames on one slide ---
% Slightly smaller body text + tighter lists/blocks. Frame titles and the
% Madrid chrome keep their normal size, so the deck still reads as Madrid,
% just with more vertical room for content.
\setbeamerfont{normal text}{size=\small}
\setbeamertemplate{itemize/enumerate body begin}{\small}
\setbeamertemplate{itemize/enumerate subbody begin}{\footnotesize}
% Tighter block spacing.
\setbeamertemplate{block begin}{%
  \par\vskip2pt%
  \begin{beamercolorbox}[colsep*=.75ex,leftskip=1ex,rightskip=1ex]{block title}%
    \usebeamerfont*{block title}\insertblocktitle%
  \end{beamercolorbox}%
  {\parskip0pt\vskip-1pt}%
  \begin{beamercolorbox}[colsep*=.75ex,vmode,leftskip=1ex,rightskip=1ex]{block body}%
  \vskip1pt}
\setbeamertemplate{block end}{\end{beamercolorbox}\vskip2pt}

\usepackage{amsmath, amssymb}
\usepackage{booktabs}
\usepackage{xcolor}
\usepackage{listings}
\usepackage[skins, breakable]{tcolorbox}

% --- Book colour palette (kept in sync with book/preamble.tex) ---
\definecolor{linkblue}{RGB}{14, 75, 140}
\definecolor{deepdivebg}{RGB}{235, 244, 255}
\definecolor{narrativebg}{RGB}{255, 252, 240}
\definecolor{narrativerule}{RGB}{180, 130, 40}
\definecolor{steelblue}{RGB}{70, 130, 180}
\definecolor{forestgreen}{RGB}{34, 139, 34}
\definecolor{crimson}{RGB}{220, 20, 60}
\definecolor{darkorange}{RGB}{255, 140, 0}
\definecolor{codebg}{RGB}{248, 248, 248}
\definecolor{coderule}{RGB}{210, 210, 210}
\definecolor{keywordblue}{RGB}{0, 100, 180}
\definecolor{commentgray}{RGB}{120, 120, 120}
\definecolor{stringgreen}{RGB}{30, 130, 30}

% Tie the Madrid structural colour to the book's link blue.
\setbeamercolor{structure}{fg=linkblue}
\setbeamercolor{title}{fg=white}
\setbeamercolor{frametitle}{fg=white}

% --- "the bigger picture" : narrative / intuition framing -----------
\newtcolorbox{bigpicture}{%
  enhanced, breakable, boxsep=2pt,
  colback  = narrativebg,
  colframe = narrativerule,
  boxrule  = 0pt, toprule = 1.4pt, bottomrule = 0.4pt,
  leftrule = 0pt, rightrule = 0pt, arc = 0pt,
  left = 6pt, right = 6pt, top = 4pt, bottom = 5pt,
  fonttitle = \footnotesize\scshape\color{narrativerule},
  title = {the bigger picture}, attach title to upper = {\par\smallskip},
}

% --- "under the hood" : the mechanism / technical detail ------------
\newtcolorbox{underhood}{%
  enhanced, breakable, boxsep=2pt,
  colback  = deepdivebg,
  colframe = linkblue,
  boxrule  = 0pt, toprule = 1.4pt, bottomrule = 0.4pt,
  leftrule = 0pt, rightrule = 0pt, arc = 0pt,
  left = 6pt, right = 6pt, top = 4pt, bottom = 5pt,
  fonttitle = \footnotesize\scshape\color{linkblue},
  title = {under the hood}, attach title to upper = {\par\smallskip},
}

% --- Python code style ---------------------------------------------
\lstset{
  basicstyle    = \ttfamily\footnotesize,
  keywordstyle  = \color{keywordblue}\bfseries,
  commentstyle  = \color{commentgray}\itshape,
  stringstyle   = \color{stringgreen},
  showstringspaces = false,
  language      = Python,
  breaklines    = true,
  backgroundcolor = \color{codebg},
  frame         = single,
  rulecolor     = \color{coderule},
  numbers       = none,
  aboveskip     = 4pt, belowskip = 4pt,
}

% --- Appendix conventions ------------------------------------------
% \appendixstart{Title} : begins the "extra material" appendix. Frames
% after it are not counted toward the main deck's frame total, keeping
% the core talk lean while preserving the deeper / optional content.
\newcommand{\appendixstart}[1]{%
  \appendix
  \setbeamertemplate{frametitle continuation}{}%
  \begin{frame}[noframenumbering]
    \begin{center}
      {\Large\scshape\color{linkblue} Appendix}\\[4pt]
      {\large #1}\\[10pt]
      {\footnotesize Extra / optional material — beyond the core lecture.}
    \end{center}
  \end{frame}
}
% Use \begin{frame}[noframenumbering]{...} for each appendix frame.

%   \title{...}\subtitle{...}\author{...}\institute{...}\date{\today}
%   \begin{document} ... \end{document}
% =====================================================================

\usetheme{Madrid}
\usecolortheme{whale}
\setbeamertemplate{navigation symbols}{}
\setbeamertemplate{caption}[numbered]
\setbeamercovered{transparent}

% --- Density controls: keep dense, equation-heavy frames on one slide ---
% Slightly smaller body text + tighter lists/blocks. Frame titles and the
% Madrid chrome keep their normal size, so the deck still reads as Madrid,
% just with more vertical room for content.
\setbeamerfont{normal text}{size=\small}
\setbeamertemplate{itemize/enumerate body begin}{\small}
\setbeamertemplate{itemize/enumerate subbody begin}{\footnotesize}
% Tighter block spacing.
\setbeamertemplate{block begin}{%
  \par\vskip2pt%
  \begin{beamercolorbox}[colsep*=.75ex,leftskip=1ex,rightskip=1ex]{block title}%
    \usebeamerfont*{block title}\insertblocktitle%
  \end{beamercolorbox}%
  {\parskip0pt\vskip-1pt}%
  \begin{beamercolorbox}[colsep*=.75ex,vmode,leftskip=1ex,rightskip=1ex]{block body}%
  \vskip1pt}
\setbeamertemplate{block end}{\end{beamercolorbox}\vskip2pt}

\usepackage{amsmath, amssymb}
\usepackage{booktabs}
\usepackage{xcolor}
\usepackage{listings}
\usepackage[skins, breakable]{tcolorbox}

% --- Book colour palette (kept in sync with book/preamble.tex) ---
\definecolor{linkblue}{RGB}{14, 75, 140}
\definecolor{deepdivebg}{RGB}{235, 244, 255}
\definecolor{narrativebg}{RGB}{255, 252, 240}
\definecolor{narrativerule}{RGB}{180, 130, 40}
\definecolor{steelblue}{RGB}{70, 130, 180}
\definecolor{forestgreen}{RGB}{34, 139, 34}
\definecolor{crimson}{RGB}{220, 20, 60}
\definecolor{darkorange}{RGB}{255, 140, 0}
\definecolor{codebg}{RGB}{248, 248, 248}
\definecolor{coderule}{RGB}{210, 210, 210}
\definecolor{keywordblue}{RGB}{0, 100, 180}
\definecolor{commentgray}{RGB}{120, 120, 120}
\definecolor{stringgreen}{RGB}{30, 130, 30}

% Tie the Madrid structural colour to the book's link blue.
\setbeamercolor{structure}{fg=linkblue}
\setbeamercolor{title}{fg=white}
\setbeamercolor{frametitle}{fg=white}

% --- "the bigger picture" : narrative / intuition framing -----------
\newtcolorbox{bigpicture}{%
  enhanced, breakable, boxsep=2pt,
  colback  = narrativebg,
  colframe = narrativerule,
  boxrule  = 0pt, toprule = 1.4pt, bottomrule = 0.4pt,
  leftrule = 0pt, rightrule = 0pt, arc = 0pt,
  left = 6pt, right = 6pt, top = 4pt, bottom = 5pt,
  fonttitle = \footnotesize\scshape\color{narrativerule},
  title = {the bigger picture}, attach title to upper = {\par\smallskip},
}

% --- "under the hood" : the mechanism / technical detail ------------
\newtcolorbox{underhood}{%
  enhanced, breakable, boxsep=2pt,
  colback  = deepdivebg,
  colframe = linkblue,
  boxrule  = 0pt, toprule = 1.4pt, bottomrule = 0.4pt,
  leftrule = 0pt, rightrule = 0pt, arc = 0pt,
  left = 6pt, right = 6pt, top = 4pt, bottom = 5pt,
  fonttitle = \footnotesize\scshape\color{linkblue},
  title = {under the hood}, attach title to upper = {\par\smallskip},
}

% --- Python code style ---------------------------------------------
\lstset{
  basicstyle    = \ttfamily\footnotesize,
  keywordstyle  = \color{keywordblue}\bfseries,
  commentstyle  = \color{commentgray}\itshape,
  stringstyle   = \color{stringgreen},
  showstringspaces = false,
  language      = Python,
  breaklines    = true,
  backgroundcolor = \color{codebg},
  frame         = single,
  rulecolor     = \color{coderule},
  numbers       = none,
  aboveskip     = 4pt, belowskip = 4pt,
}

% --- Appendix conventions ------------------------------------------
% \appendixstart{Title} : begins the "extra material" appendix. Frames
% after it are not counted toward the main deck's frame total, keeping
% the core talk lean while preserving the deeper / optional content.
\newcommand{\appendixstart}[1]{%
  \appendix
  \setbeamertemplate{frametitle continuation}{}%
  \begin{frame}[noframenumbering]
    \begin{center}
      {\Large\scshape\color{linkblue} Appendix}\\[4pt]
      {\large #1}\\[10pt]
      {\footnotesize Extra / optional material — beyond the core lecture.}
    \end{center}
  \end{frame}
}
% Use \begin{frame}[noframenumbering]{...} for each appendix frame.

%   \title{...}\subtitle{...}\author{...}\institute{...}\date{\today}
%   \begin{document} ... \end{document}
% =====================================================================

\usetheme{Madrid}
\usecolortheme{whale}
\setbeamertemplate{navigation symbols}{}
\setbeamertemplate{caption}[numbered]
\setbeamercovered{transparent}

% --- Density controls: keep dense, equation-heavy frames on one slide ---
% Slightly smaller body text + tighter lists/blocks. Frame titles and the
% Madrid chrome keep their normal size, so the deck still reads as Madrid,
% just with more vertical room for content.
\setbeamerfont{normal text}{size=\small}
\setbeamertemplate{itemize/enumerate body begin}{\small}
\setbeamertemplate{itemize/enumerate subbody begin}{\footnotesize}
% Tighter block spacing.
\setbeamertemplate{block begin}{%
  \par\vskip2pt%
  \begin{beamercolorbox}[colsep*=.75ex,leftskip=1ex,rightskip=1ex]{block title}%
    \usebeamerfont*{block title}\insertblocktitle%
  \end{beamercolorbox}%
  {\parskip0pt\vskip-1pt}%
  \begin{beamercolorbox}[colsep*=.75ex,vmode,leftskip=1ex,rightskip=1ex]{block body}%
  \vskip1pt}
\setbeamertemplate{block end}{\end{beamercolorbox}\vskip2pt}

\usepackage{amsmath, amssymb}
\usepackage{booktabs}
\usepackage{xcolor}
\usepackage{listings}
\usepackage[skins, breakable]{tcolorbox}

% --- Book colour palette (kept in sync with book/preamble.tex) ---
\definecolor{linkblue}{RGB}{14, 75, 140}
\definecolor{deepdivebg}{RGB}{235, 244, 255}
\definecolor{narrativebg}{RGB}{255, 252, 240}
\definecolor{narrativerule}{RGB}{180, 130, 40}
\definecolor{steelblue}{RGB}{70, 130, 180}
\definecolor{forestgreen}{RGB}{34, 139, 34}
\definecolor{crimson}{RGB}{220, 20, 60}
\definecolor{darkorange}{RGB}{255, 140, 0}
\definecolor{codebg}{RGB}{248, 248, 248}
\definecolor{coderule}{RGB}{210, 210, 210}
\definecolor{keywordblue}{RGB}{0, 100, 180}
\definecolor{commentgray}{RGB}{120, 120, 120}
\definecolor{stringgreen}{RGB}{30, 130, 30}

% Tie the Madrid structural colour to the book's link blue.
\setbeamercolor{structure}{fg=linkblue}
\setbeamercolor{title}{fg=white}
\setbeamercolor{frametitle}{fg=white}

% --- "the bigger picture" : narrative / intuition framing -----------
\newtcolorbox{bigpicture}{%
  enhanced, breakable, boxsep=2pt,
  colback  = narrativebg,
  colframe = narrativerule,
  boxrule  = 0pt, toprule = 1.4pt, bottomrule = 0.4pt,
  leftrule = 0pt, rightrule = 0pt, arc = 0pt,
  left = 6pt, right = 6pt, top = 4pt, bottom = 5pt,
  fonttitle = \footnotesize\scshape\color{narrativerule},
  title = {the bigger picture}, attach title to upper = {\par\smallskip},
}

% --- "under the hood" : the mechanism / technical detail ------------
\newtcolorbox{underhood}{%
  enhanced, breakable, boxsep=2pt,
  colback  = deepdivebg,
  colframe = linkblue,
  boxrule  = 0pt, toprule = 1.4pt, bottomrule = 0.4pt,
  leftrule = 0pt, rightrule = 0pt, arc = 0pt,
  left = 6pt, right = 6pt, top = 4pt, bottom = 5pt,
  fonttitle = \footnotesize\scshape\color{linkblue},
  title = {under the hood}, attach title to upper = {\par\smallskip},
}

% --- Python code style ---------------------------------------------
\lstset{
  basicstyle    = \ttfamily\footnotesize,
  keywordstyle  = \color{keywordblue}\bfseries,
  commentstyle  = \color{commentgray}\itshape,
  stringstyle   = \color{stringgreen},
  showstringspaces = false,
  language      = Python,
  breaklines    = true,
  backgroundcolor = \color{codebg},
  frame         = single,
  rulecolor     = \color{coderule},
  numbers       = none,
  aboveskip     = 4pt, belowskip = 4pt,
}

% --- Appendix conventions ------------------------------------------
% \appendixstart{Title} : begins the "extra material" appendix. Frames
% after it are not counted toward the main deck's frame total, keeping
% the core talk lean while preserving the deeper / optional content.
\newcommand{\appendixstart}[1]{%
  \appendix
  \setbeamertemplate{frametitle continuation}{}%
  \begin{frame}[noframenumbering]
    \begin{center}
      {\Large\scshape\color{linkblue} Appendix}\\[4pt]
      {\large #1}\\[10pt]
      {\footnotesize Extra / optional material — beyond the core lecture.}
    \end{center}
  \end{frame}
}
% Use \begin{frame}[noframenumbering]{...} for each appendix frame.


\title{Practical Session 12: Explaining Financial LLM Decisions}
\subtitle{Lecture 12 --- Hands-On: SHAP by Hand, Faithfulness, and an ECOA Pipeline}
\author{Juan F.\ Imbet}
\institute{Paris Dauphine -- PSL University}
\date{\today}

\begin{document}

\begin{frame}\titlepage\end{frame}

% ============================================================
% FRAME 1: Session Overview
% ============================================================
\begin{frame}{Session Overview}
\textbf{Duration:} 1 hour (50 min problems + 10 min debrief)

\medskip
\textbf{Agenda:}
\begin{enumerate}
  \item \textbf{Recap} (5 min): Shapley formula, faithfulness metrics, the
        constrained-generation pipeline.
  \item \textbf{Problem 1} (12 min): Compute Shapley values by hand for a 2-feature
        credit game; verify the efficiency property.
  \item \textbf{Problem 2} (12 min): Score sufficiency \& comprehensiveness; judge
        whether an explanation is faithful.
  \item \textbf{Problem 3} (8 min): Audit an LLM-drafted denial letter for ECOA
        compliance.
  \item \textbf{Case Study} (15 min): Run the deterministic SHAP attribution
        generator in \texttt{code/notebooks/12-xai-explainability/}; vary inputs.
\end{enumerate}

\medskip
\begin{block}{Discussion (8 min)}
When is a \emph{plausible} explanation worse than no explanation at all?
\end{block}
\end{frame}

% ============================================================
% FRAME 2: Recap — Shapley
% ============================================================
\begin{frame}{Recap: Shapley Values and Efficiency}
You will need this for Problem 1.

\begin{underhood}
\[
\phi_i(v) = \sum_{S \subseteq N \setminus \{i\}}
\frac{|S|!\,(n-|S|-1)!}{n!}\bigl[v(S \cup \{i\}) - v(S)\bigr].
\]
For $n=2$, each feature has exactly two orderings, so
\[
\phi_1 = \tfrac{1}{2}\bigl[v(\{1\})-v(\emptyset)\bigr]
       + \tfrac{1}{2}\bigl[v(\{1,2\})-v(\{2\})\bigr].
\]
\end{underhood}

\medskip
\textbf{Efficiency (the check):}
\[
\phi_1 + \phi_2 = v(\{1,2\}) - v(\emptyset) = f(x) - \mathbb{E}[f(X)].
\]
The attributions must sum to the gap between prediction and baseline.
\end{frame}

% ============================================================
% FRAME 3: Recap — Faithfulness + pipeline
% ============================================================
\begin{frame}{Recap: Faithfulness Metrics and the Safe Pipeline}
\textbf{Sufficiency} (keep only the explained set $e$):
\[
\text{Suff}(e,f,x) = 1 - |f(x) - f(x_{e(x)})| \quad (\to 1 \text{ is best}).
\]
\textbf{Comprehensiveness} (remove the explained set):
\[
\text{Comp}(e,f,x) = |f(x) - f(x_{\bar{e}(x)})| \quad (\text{high is good}).
\]

\medskip
\begin{block}{The constrained explanation pipeline (for Problem 3)}
Decide $\to$ attribute (SHAP/IG) $\to$ \emph{constrained} LLM translation $\to$
compliance filter $\to$ human review. The generator may refer \textbf{only} to the
attribution feed; it must not invent factors or reference protected
characteristics (ECOA).
\end{block}
\end{frame}

% ============================================================
% FRAME 4: Problem 1
% ============================================================
\begin{frame}{Problem 1: Shapley Values for a Credit Game}
A credit model scores an applicant with two features:
$x_1=$ debt-to-income, $x_2=$ revenue trend. The value function $v(S)=$ model's
$P(\text{creditworthy})$ when only the features in $S$ are known:

\medskip
\begin{tabular}{lc}
\toprule
Coalition $S$ & $v(S)$ \\
\midrule
$\emptyset$ (baseline) & $0.50$ \\
$\{1\}$ (DTI only)     & $0.32$ \\
$\{2\}$ (revenue only) & $0.41$ \\
$\{1,2\}$ (both)       & $0.14$ \\
\bottomrule
\end{tabular}

\medskip
\textbf{Task A.} Compute $\phi_1$ and $\phi_2$ from the $n=2$ formula.\\
\textbf{Task B.} Verify the efficiency property $\phi_1+\phi_2 = v(\{1,2\})-v(\emptyset)$.\\
\textbf{Task C.} Which feature is the stronger driver of the rejection? In an
ECOA notice, which reason is listed first?
\end{frame}

% ============================================================
% FRAME 5: Problem 2
% ============================================================
\begin{frame}{Problem 2: Is This Explanation Faithful?}
A text classifier outputs $f(x)=0.14$ for a loan application (baseline $0.50$).
An XAI method claims the explained set $e(x) = \{$``struggled'', ``declining'',
``down 30\%''$\}$ accounts for the decision. You run two ablations:

\medskip
\begin{itemize}
  \item Keep \emph{only} $e(x)$ (mask all other tokens): $f(x_{e(x)}) = 0.19$.
  \item Remove \emph{only} $e(x)$ (keep the rest): $f(x_{\bar{e}(x)}) = 0.44$.
\end{itemize}

\medskip
\textbf{Task A.} Compute the sufficiency $\text{Suff}(e,f,x)$.\\
\textbf{Task B.} Compute the comprehensiveness $\text{Comp}(e,f,x)$.\\
\textbf{Task C.} Is this explanation \emph{faithful}? Interpret both numbers ---
what would a low sufficiency \emph{and} a low comprehensiveness have told you?\\
\textbf{Task D.} The attention map for the same input highlights ``restaurant''
and ``market'' instead. After ablating those two tokens, $f$ is unchanged at
$0.14$. What does that imply about the attention explanation?
\end{frame}

% ============================================================
% FRAME 6: Problem 3
% ============================================================
\begin{frame}{Problem 3: Audit an LLM-Drafted Denial Letter}
An LLM was given this attribution feed (DTI $-0.18$, revenue $-0.12$, cash
reserve $-0.07$) and produced the draft below. \textbf{Find every compliance
defect} against ECOA / Reg.\ B.

\medskip
\begin{exampleblock}{Draft denial letter}
``Dear Applicant, your application was declined. Your overall credit profile did
not meet our criteria. We also note that applicants in your area and of your
demographic tend to default more often, and your business outlook seemed weak.
We wish you the best.''
\end{exampleblock}

\medskip
\textbf{Task.} Identify which of the following the draft violates and why:
\begin{enumerate}\small
  \item Vagueness (specificity requirement).
  \item Introduction of a factor \emph{not} in the attribution feed.
  \item Reference to a protected characteristic.
  \item Missing required disclosures (statement-of-reasons / consumer-report rights).
\end{enumerate}
Rewrite the first sentence to be \emph{specific and compliant} using the feed.
\end{frame}

% ============================================================
% FRAME 7: Case Study — Intro
% ============================================================
\begin{frame}{Case Study: Deterministic SHAP Attribution Figure}
\textbf{Goal:} reproduce the token-level SHAP attribution chart from the chapter's
credit example, then modify it.

\medskip
\textbf{Code:} \texttt{code/notebooks/12-xai-explainability/gen\_shap\_attribution.py}

\medskip
\textbf{Why deterministic?} The script hard-codes the documented attribution
values from the example --- it depends on \emph{no} live model or data, so every
student reproduces an identical figure. This is exactly the
\emph{reproducibility} property auditors require.

\medskip
\begin{alertblock}{Prerequisites}
\texttt{pip install matplotlib}\\
Run: \texttt{python3 gen\_shap\_attribution.py} (writes a PDF into the book's
\texttt{figures/} directory and prints the sum of shown tokens).
\end{alertblock}
\end{frame}

% ============================================================
% FRAME 8: Case Study — Code
% ============================================================
\begin{frame}[fragile]{Case Study: The Code}
\begin{lstlisting}[basicstyle=\ttfamily\scriptsize]
TOKENS = [
    ("struggled",   -0.089),
    ("declining",   -0.073),
    ("down 30%",    -0.061),
    ("competitive", -0.047),
    ("restaurant",  -0.038),
    ("pandemic",    -0.029),
    ("operates",    +0.012),
]

# sort most-negative at the bottom for a waterfall read
items  = sorted(TOKENS, key=lambda t: t[1])
labels = [t[0] for t in items]
values = [t[1] for t in items]
colors = ["#d62728" if v < 0 else "#2ca02c" for v in values]

fig, ax = plt.subplots(figsize=(7.2, 4.2))
ax.barh(range(len(values)), values, color=colors)
ax.axvline(0.0, color="black", linewidth=0.8)
ax.set_xlabel("SHAP value (contribution to P(creditworthy))")
fig.savefig(OUT, bbox_inches="tight")
print("sum of shown tokens =", round(sum(values), 3))
\end{lstlisting}
\textbf{Baseline} $0.50 \to$ model output $0.14$; total gap $= -0.36$.
\end{frame}

% ============================================================
% FRAME 9: Case Study — Questions
% ============================================================
\begin{frame}{Case Study: Diagnostic Questions}
Run the script, then answer:

\medskip
\textbf{Q1.} What is the printed sum of the seven shown tokens? Compare it to the
full prediction gap $0.14 - 0.50 = -0.36$. Why is the sum of the \emph{shown}
tokens not exactly $-0.36$? (Hint: \emph{completeness}.)

\medskip
\textbf{Q2.} Compute the SHAP \emph{completeness} ratio
$\sum_{i\in e} \phi_i / (f(x)-\mathbb{E}[f(X)])$ for the seven tokens. What
fraction of the decision do they explain?

\medskip
\textbf{Q3.} Change ``operates'' from $+0.012$ to $-0.030$. Does the colour of its
bar change? What would such a sign flip imply for the adverse-action narrative?

\medskip
\textbf{Q4.} The figure is \emph{deterministic}. Contrast this with running
\emph{KernelSHAP} or \emph{LIME} live on the model. Which property required for
audit does the deterministic version guarantee that a live, unseeded run does not?
\end{frame}

% ============================================================
% FRAME 10: Discussion
% ============================================================
\begin{frame}{Discussion: When Is a Plausible Explanation Dangerous?}
For each scenario, decide whether the explanation is safe to ship, and what extra
check you would require.

\medskip
\begin{enumerate}
  \item A retail-facing chatbot generates a fluent paragraph explaining a loan
        rejection. It was \emph{not} given the upstream SHAP feed.
  \item An attention heat-map highlights tokens an expert agrees are relevant, but
        ablating them does not change the prediction.
  \item A robo-adviser's suitability narrative cites a 9\% expected return that does
        not appear in the product's Key Information Document.
\end{enumerate}

\medskip
\begin{block}{Frame your answer around:}
faithfulness vs.\ plausibility, the separation-of-concerns pipeline, protected
characteristics, and reproducibility for audit.
\end{block}
\end{frame}

% ============================================================
% FRAME 11: Solution — Problem 1
% ============================================================
\begin{frame}{Solution: Problem 1 --- Shapley Values}
\textbf{Task A.}
\[
\phi_1 = \tfrac{1}{2}[v(\{1\})-v(\emptyset)] + \tfrac{1}{2}[v(\{1,2\})-v(\{2\})]
       = \tfrac{1}{2}(0.32-0.50) + \tfrac{1}{2}(0.14-0.41)
\]
\[
= \tfrac{1}{2}(-0.18) + \tfrac{1}{2}(-0.27) = -0.090 - 0.135 = \mathbf{-0.225}.
\]
\[
\phi_2 = \tfrac{1}{2}(0.41-0.50) + \tfrac{1}{2}(0.14-0.32)
       = \tfrac{1}{2}(-0.09) + \tfrac{1}{2}(-0.18) = \mathbf{-0.135}.
\]

\textbf{Task B (efficiency).}
$\phi_1 + \phi_2 = -0.225 - 0.135 = -0.360 = v(\{1,2\}) - v(\emptyset) = 0.14-0.50.$ \checkmark

\medskip
\textbf{Task C.} $|\phi_1| > |\phi_2|$: \textbf{debt-to-income} is the stronger
driver $\Rightarrow$ listed first in the ECOA notice.
\end{frame}

% ============================================================
% FRAME 12: Solution — Problem 2
% ============================================================
\begin{frame}{Solution: Problem 2 --- Faithfulness}
\textbf{Task A (sufficiency).}
$\text{Suff} = 1 - |0.14 - 0.19| = 1 - 0.05 = \mathbf{0.95}$ --- high: keeping
only $e(x)$ nearly reproduces the prediction.

\medskip
\textbf{Task B (comprehensiveness).}
$\text{Comp} = |0.14 - 0.44| = \mathbf{0.30}$ --- large: removing $e(x)$ moves the
prediction most of the way back toward the $0.50$ baseline.

\medskip
\textbf{Task C.} High sufficiency \emph{and} high comprehensiveness $\Rightarrow$
the explanation is \textbf{faithful}: the three tokens both suffice for and are
necessary to the decision. Low sufficiency would mean $e(x)$ omits real drivers;
low comprehensiveness would mean the model routes around $e(x)$.

\medskip
\textbf{Task D.} The attention explanation \textbf{fails comprehensiveness}
(ablating its tokens leaves $f=0.14$ unchanged). It is \emph{plausible} but not
\emph{faithful} --- not audit-grade.
\end{frame}

% ============================================================
% FRAME 13: Solution — Problem 3
% ============================================================
\begin{frame}{Solution: Problem 3 --- Denial-Letter Audit}
\textbf{Defects in the draft:}
\begin{enumerate}\small
  \item \textbf{Vagueness:} ``credit profile did not meet our criteria'' --- no
        specific, quantified reasons. Violates Reg.\ B specificity.
  \item \textbf{Unlisted factor:} ``business outlook seemed weak'' is not in the
        feed; the generator hallucinated it.
  \item \textbf{Protected characteristic:} ``applicants \ldots\ of your
        demographic'' --- prohibited basis under ECOA. Severe.
  \item \textbf{Missing disclosures:} no statement-of-reasons right, no
        free-consumer-report-copy notice.
\end{enumerate}

\medskip
\textbf{Compliant rewrite of the first sentence:}
\begin{exampleblock}{}
``Your application was declined primarily because your debt-to-income ratio of
54\% exceeded our maximum of 45\%, your revenue declined 28\% year-on-year, and
your cash reserves covered under one month of operating expenses.''
\end{exampleblock}
Specific, quantified, drawn \emph{only} from the feed, no protected factors.
\end{frame}

% ============================================================
% FRAME 14: Wrap-up
% ============================================================
\begin{frame}{Wrap-Up: Key Takeaways}
\textbf{From Problem 1:} Shapley values are marginal contributions averaged over
orderings; the efficiency property is your built-in correctness check, and
$|\phi_i|$ orders the ECOA reasons.

\medskip
\textbf{From Problem 2:} Faithfulness $=$ sufficiency \emph{and} comprehensiveness.
Attention can be plausible yet fail comprehensiveness --- not audit-grade.

\medskip
\textbf{From Problem 3:} The dangerous failures are hallucinated factors and
protected-characteristic leakage. Constrain the generator + add a compliance
filter + human review.

\medskip
\textbf{From the Case Study:} Deterministic, seeded explanations are reproducible
--- the property audits demand and live LIME/KernelSHAP runs do not guarantee.

\medskip
\begin{block}{Next practical}
Privacy and local models: running explainable pipelines on-premises when client
data cannot leave the institution.
\end{block}
\end{frame}

% ============================================================
% APPENDIX
% ============================================================
\appendixstart{Stretch Problems}

\begin{frame}[noframenumbering]{Appendix: Stretch Problem --- Counterfactual Search}
A scoring model uses two structured features, DTI $d$ and cash-reserve months $c$.
The decision boundary (approve iff score $\geq 0.5$) is
$\text{score} = 0.5 - 1.2(d - 0.45) + 0.1(c - 3)$, where $d$ is a fraction.

\medskip
An applicant has $d=0.54$, $c=0.8$.
\begin{enumerate}\small
  \item Compute the current score. Is the applicant approved?
  \item Find the minimal \emph{single-feature} change in $d$ (holding $c$ fixed)
        that flips the decision. Express it as an actionable counterfactual.
  \item Repeat for $c$. Which counterfactual is the smaller \emph{relative} change?
  \item Why might the $d$-counterfactual be preferred even if the $c$-change is
        numerically smaller? (Hint: feasibility \& cost-reflecting distance.)
\end{enumerate}

\medskip
\textbf{Sketch:} score $= 0.5 - 1.2(0.09) + 0.1(-2.2) = 0.5 - 0.108 - 0.22 =
0.172 < 0.5 \Rightarrow$ declined. Approval needs the linear term $\geq 0$: e.g.\
$d \leq 0.45 - (0.1(c-3))/1.2$. Articulate the resulting threshold as a plain
counterfactual.
\end{frame}

\end{document}
